% Standalone appendix document - compiles on its own:
%   latexmk -pdf 02-method-and-process.tex
\documentclass[a4paper,11pt]{article}
\usepackage[english]{babel}
\usepackage{../../appendix-style}
\usepackage{pgfplotstable}
\pgfplotsset{compat=1.18}

\hypersetup{
    colorlinks=true,
    linkcolor=blue,
    urlcolor=blue
}

\addbibresource{../../../assets/references.bib}

\title{10.2 Method and Process}
\author{SW4PRJ4 -- Group 3}
\date{\today}

\begin{document}
\maketitle

% Content goes here...
\begin{itemize}
    \item Frederik volunteered to be Scrummaster. Everyone was in agreement.
    \item Product owner is a blue elephant named Bodil.
    \item Sprint planning is done in Kaneo. Different group members had different experiences with Trello and Jira, so we started in Notion and moved to Kaneo on 9/9.
\end{itemize}

\subsection*{2/9}
\begin{itemize}
    \item We went with User Stories to try something else than fully dressed use cases. User Stories appealed to the team since this is as close to the consumer as possible, expressed in the customer's own words and language, which made it easier to validate requirements and keep the focus on delivering what the consumer wants.
    \item Requirement-specification $\rightarrow$ A lot of thought went into the database. What relations should be between each ``member'' of the database.
    \item Created user stories before the beginning of sprint\#1. We followed the setup from brightspace.
    \item During requirement specifications the group spend a lot of time on discussing the scope of the project and how the database should be constructed. Should every customer be able to swap a single food item or should it be x amount of packets which content is irreplaceable? We chose certain packets with irreplaceable content since we presumed that the vast majority of customers would pick either a ``standard package'', a ``vegetarian package'' or ``gluten-free'' package. Whole household = one package. One package per person was probably feasible as well, but we chose to get with it as a ``Could'' in MoSCoW since we deemed it a minor modification to the system.
    \item Concept clarification started.
\end{itemize}

\subsection*{9/9}
\begin{itemize}
    \item Fred presents how the project could be shaped using AI as development-tool. Unanimous agreement. See \href{run:./Arkitektur-ide.pdf}{Arkitektur-idé}.
    \item Moreover, Notion was ditched and Kaneo was chosen. Kaneo is an open-source alternative to Trello. It is easy to use and has what we need, without the information overload of larger tools such as Jira. It is also compatible with Claude. This way, Claude will be able to create and chose tickets and complete these.
\end{itemize}

\subsection*{15/9}
\begin{itemize}
    \item Review \#1 with group 4. Their feedback led to FR/NFR numbering and testable NFRs in the requirements specification.
    \item We agreed that step 2 is to add an external authority for login, and discussed Firebase, Keycloak, Clerk and OAuth.
    \item We agreed on one shared deployment instead of everyone running the system locally, because an administrator needs to work on the same data as the customers.
\end{itemize}

\subsection*{16/9}
\begin{itemize}
    \item The project skeleton was set up: the layered .NET solution, test projects, the Expo app, the React admin panel and the shared api-client in one npm workspace.
    \item Keycloak was chosen as identity provider and runs in docker compose. Coolify on our own server was chosen for deployment. The reasons are in appendix 2.1.
    \item We use GitHub instead of GitLab, which the first architecture pitch planned for. GitHub Actions builds and tests on every pull request.
    \item We chose GitHub Flow with Conventional Commits and named branches. Git hooks and a GitHub Actions check make sure the rules are followed. See appendix 10.1.
    \item We changed our mind about Claude and Kaneo: the AI may read, comment on and create tasks, but moving, assigning and deleting tasks is done by the group. Planning the sprint is part of what we have to learn ourselves.
    \item We decided that all text in the apps is in English.
\end{itemize}

\subsection*{23/9}
\begin{itemize}
    \item Sprints are two weeks, from Wednesday to Tuesday. Sprint 1 runs from 16/9 to 29/9. The sprint calendar is in appendix 10.1.
    \item The Backlog view in Kaneo is our product backlog, and the board is the sprint backlog.
    \item Tasks get labels for story points, sprint and epic, so we can measure velocity and trace a task back to the requirements.
\end{itemize}

\section*{Velocity}
Committed is the sum of story points in \emph{To Do} at sprint planning. Done is the sum in \emph{Done} at sprint review. Spillover is the points that move on to the next sprint. Backlog size is the sum of all estimated points in the project at sprint review.

The numbers are kept in \texttt{docs/assets/data/sprint-velocity.csv}. The table below and the burn-up chart in chapter~2 of the report are both generated from that file, so it is the only place the numbers are entered.

\begin{center}
\pgfplotstabletypeset[
    col sep=comma,
    empty cells with={--},
    every head row/.style={after row=\hline},
    columns/sprint/.style={column name=Sprint},
    columns/committed/.style={column name=Committed},
    columns/done/.style={column name=Done},
    columns/spillover/.style={column name=Spillover},
    columns/scope/.style={column name=Backlog size},
]{../../../assets/data/sprint-velocity.csv}
\end{center}

\section*{Sprint Goals}
\begin{center}
\begin{tabular}{@{}llp{9cm}@{}}
Sprint & Period & Sprint goal \\ \hline
1 & 16/9 -- 29/9 & \\
2 & 30/9 -- 13/10 & \\
3 & 14/10 -- 27/10 & \\
4 & 28/10 -- 10/11 & \\
5 & 11/11 -- 24/11 & \\
6 & 25/11 -- 8/12 & \\
\end{tabular}
\end{center}

\printbibliography
\end{document}
